\chapter{Forme Differenziali} 

\section{forme differenziali lineari}

Se $f\colon \Omega\subset \RR^n \to \RR$ è una funzione differenziabile il suo 
differenziale $df \colon \Omega \to \mathcal L (\RR^n, \RR)$ è una forma differenziale lineare
($1$-forma).
Si tratta di una funzione che ad ogni punto $x\in \Omega$ associa una forma lineare 
$df(x)\colon \RR^n \to \RR$ e dunque ad ogni $v\in \RR^n$ associa un numero reale 
$df(x)[v]$. 
Se sfruttiamo il prodotto scalare standard su $\RR^n$ 
possiamo identificare l'applicazione lineare $df(x)$ con il vettore $\nabla f(x)$:
\[
    df(x)[v] = \frac{\partial f}{\partial v}(x) = \nabla f(x) \cdot v.
\]
Infatti il prodotto scalare su uno spazio vettoriale $V$ 
induce un isomorfismo $V \to V^*$, $v \mapsto (w \mapsto v \cdot w)$.

\begin{definition}[forme differenziali lineari]
    Sia $\Omega\subset \RR^n$ un aperto.
    Denotiamo con $\Lambda^1(\Omega)$ l'insieme di tutte le funzioni $\omega\colon \Omega \to (\RR^n)^*$
    che ad ogni punto $x\in \Omega$ associano una forma lineare $\omega(x)\colon \RR^n \to \RR$.
\end{definition}

Grazie all'isomorfismo dato dal prodotto scalare sappiamo 
che la dimensione di $(\RR^n)^*$ è $n$. 
Determiniamo una base canonica di questo spazio.
Denotiamo con $e_1, \dots e_n$ la base canonica di $\RR^n$.
Denotiamo con $x_1, \dots x_n$ le funzioni $\RR^n \to \RR$ 
che ad ogni punto $p\in \RR^n$ associano la $k$-esima coordinata di $p$:
$x_k(p) = p_k = p\cdot e_k$.
Possiamo allora osservare che il differenziale $dx_k$ 
è tale che
\[
  dx_k(p)[e_j] = \frac{\partial x_k}{x_j} = \delta_{jk} 
  = \begin{cases}
    1 & j=k\\
    0 & j\neq k.
  \end{cases}  
\]
In effetti per ogni $p\in \Omega$ le applicazioni lineari 
$dx_1(p), \dots, dx_n(p)$ sono una base di $(\RR^n)^*$. 
Significa che ogni $\omega \in \Lambda^1(\Omega)$ può essere scritto 
nella forma 
\[
   \omega(p) = \xi_1 dx_1 + \xi_2 dx_2 + \dots + \xi_n dx_n.
\]
Dunque c'è un isomorfismo $\omega \mapsto \xi$ che mette in 
corrispondenza le forme differenziali lineari $\omega \in \Lambda^1(\Omega)$ 
con i campi vettoriali $\xi \colon \Omega\to \RR^n$.
In questo isomorfismo:
\[
  \omega(p)[v] = \xi(p) \cdot v.
\]

\section{integrale di una forma differenziale}

Se $u\colon [a,b]\to \Omega\subset \RR^n$ è una curva regolare e $\omega\in \Lambda^1(\Omega)$ 
è una $1$-forma differenziale, possiamo definire l'integrale della forma differenziale lungo la curva:
\[
 \int_u \omega \defeq = \int_a^b \omega(u(t))[u'(t)]\, dt.
\]
Osserviamo che se $\omeg$ si rappresenta con il campo vettoriale $\xi$, si ha $\omega(x)[v] = \xi(x)\cdot v$ 
e dunque 
\[
 \int_u \omega = \int_a^b \xi(u(t))\cdot u'(t)\, dt.
\]
Questo integrale si chiama \emph{circuitazione} di $\xi$ lungo $u$.
Osserviamo che questo integrale è indipendente dalla parametrizzazione della curva $u$
(se si mantiene l'orientazione)
in quanto posto $t=\phi(s)$ si ha $(u\circ \phi)'(s)\, ds = u'(\phi(s)) \phi'(s)\, ds = u'(t)\, dt$
e se $\phi\colon [A,B]\to [a,b]$ con $\phi(A)=a$ e $\phi(B)=b$ 
la formula del cambio di variabile garantisce l'uguaglianza
\[
\int_(u\circ \phi) \omega = 
\int_A^B \omega(u(\phi(s)))[(u\circ \phi)'(s)]\, ds
= \int_A^B \omega(u(\phi(s)))[u'(\phi(s))] \phi'(s)\, ds
= \int_a^b \omega(t)[u'(t)]\, dt = \int_u \omega.
\]

\section{forme chiuse, forme esatte}

Osserviamo che non tutte le forme differenziali lineari sono il differenziale
di una funzione $f\colon \Omega\to \RR$.
Infatti se $f\in C^2(\Omega)$ e $\omega = df$
allora 
\[
    \omega = \xi_1 dx_1 + \dots + \xi_n dx_n
\]
e, per il teorema di Schwartz si ha 
\begin{equation}\label{eq:forma-chiusa}
  \frac{\partial \xi_k}{\partial x_j} = \frac{\partial \xi_j}{\partial x_k}.
\end{equation}

\begin{definition}[forma chiusa]
Sia $\omega\in \Lambda^1(\Omega)$ una $1$-forma di classe $C^1$
ovvero tale che esiste $\xi\in C^1(\Omega,\RR^n)$ tale che 
$\omega(x)[v] = \xi(x)\cdot v$. 
Diremo che $\omega$ è chiusa se vale\eqref{eq:forma-chiusa}
\end{definition}

Se $df = \omega$ diremo che $f$ è una primitiva di $\omega$ 
e diremo che $\omega$ è integrabile. 

\begin{theorem}[condizione necessaria per l'integrabilità]
\label{th:forma-chiusa}
Se $\omega$ è una $1$-forma di classe $C^1$ ed esiste 
$f\in C^2(\Omega)$ tale che $\omega = df$ allora $\omega$ è chiusa. 
\end{theorem}

Per le funzioni di una variabile la codizione~\eqref{eq:forma-chiusa} è vuota, 
perché necessariamente $k=j=1$. 
E, in effetti, tutte le funzioni continue, in una variabile, sono integrabili 
per il teorema fondamentale del calcolo. 

In più variabili la condizione non è vuota. 
\begin{example}[forma differenziale non chiusa]
Sia $\omega(x,y) = y dx + dy$, $\omega \in \Lambda^1(\RR^2)$ è una forma di classe $C^\infty$ 
associata al campo vettoriale $\xi(x,y) = (y,1)$.
L'equazione~\eqref{eq:forma-chiusa} diventa $1=0$ e dunque questa forma non è chiusa e non 
può essere il differenziale di una funzione.
\end{example}

